\hypertarget{object-oriented-programming-mechanics-mro}{%
\chapter{OBJECT-ORIENTED PROGRAMMING MECHANICS \&
MRO}\label{object-oriented-programming-mechanics-mro}}

\hypertarget{method-resolution-order-mro-c3-linearization}{%
\subsection{4.1 Method Resolution Order (MRO) \& C3
Linearization}\label{method-resolution-order-mro-c3-linearization}}

In pre-2.2 versions of Python, classes were defined as ``classic
classes.'' * \textbf{Classic Classes MRO (DFLR)}: Classic class methods
were resolved using \textbf{Depth-First, Left-to-Right (DFLR)} lookup.
In a diamond inheritance hierarchy where classes share a common
ancestor, DFLR searches all the way up to the ancestor via the first
parent class, before evaluating sibling classes. This resulted in the
ancestor overriding custom methods written in sibling classes. *
\textbf{New-Style Classes MRO (C3 Linearization)}: Python 2.2 introduced
new-style classes (inheriting from \passthrough{\lstinline!object!}),
which calculate their lookup order using \textbf{C3 Linearization}. This
guarantees: 1. \textbf{Local Precedence}: Parent classes are searched in
the order declared in the class definition. 2. \textbf{Monotonicity}: A
class's MRO respects the relative ordering defined in all its parent
classes' MROs.

\hypertarget{the-c3-algorithm}{%
\subsubsection{The C3 Algorithm:}\label{the-c3-algorithm}}

Let \(L(C)\) be the MRO list of class \(C\). For a class \(C\)
inheriting from parents \(B_1, B_2, \dots, B_N\):
\[L(C) = [C] + \text{merge}(L(B_1), L(B_2), \dots, L(B_N), [B_1, B_2, \dots, B_N])\]

The \textbf{merge} operation selects a candidate head from the lists: 1.
Look at the head (first element) of the first list being merged. 2. If
this head does not appear in the tail (from index 1 to the end) of any
of the other lists being merged, it is a valid candidate. 3. Add the
candidate to \(L(C)\), remove it from all merging lists, and repeat the
merge step. 4. If the candidate head is in the tail of any other list,
skip it and check the head of the next list in the merge block. 5. If
all lists are empty, the merge completes successfully. If no candidate
head can be selected, the hierarchy is invalid (raises a
\passthrough{\lstinline!TypeError: Cannot create a consistent method resolution order!}).

\hypertarget{mathematical-step-by-step-calculation}{%
\subsubsection{Mathematical Step-by-Step
Calculation:}\label{mathematical-step-by-step-calculation}}

Let's compute the MRO for the following class definitions:

\begin{lstlisting}[language=Python]
class O(object): pass
class X(O): pass
class Y(O): pass
class A(X, Y): pass
\end{lstlisting}

We know: \[L(O) = [O, \text{object}]\]
\[L(X) = [X] + \text{merge}(L(O), [O]) = [X] + \text{merge}([O, \text{object}], [O]) = [X, O, \text{object}]\]
\[L(Y) = [Y] + \text{merge}(L(O), [O]) = [Y, O, \text{object}]\]

Now we compute \(L(A)\):
\[L(A) = [A] + \text{merge}(L(X), L(Y), [X, Y])\]
\[L(A) = [A] + \text{merge}([X, O, \text{object}], [Y, O, \text{object}], [X, Y])\]

\begin{enumerate}
\def\labelenumi{\arabic{enumi}.}
\tightlist
\item
  Evaluate head of first list: \passthrough{\lstinline!X!}. Is
  \passthrough{\lstinline!X!} in the tail of \([Y, O, \text{object}]\)
  or \([X, Y]\)? No (it is only at the head of \([X, Y]\)). We extract
  \passthrough{\lstinline!X!}:
  \[L(A) = [A, X] + \text{merge}([O, \text{object}], [Y, O, \text{object}], [Y])\]
\item
  Evaluate head of first list: \passthrough{\lstinline!O!}. Is
  \passthrough{\lstinline!O!} in the tail of \([Y, O, \text{object}]\)
  or \([Y]\)? Yes, \passthrough{\lstinline!O!} is in the tail of
  \([Y, O, \text{object}]\) (index 1). We cannot extract it.
\item
  Move to the next list \([Y, O, \text{object}]\). Evaluate head:
  \passthrough{\lstinline!Y!}. Is \passthrough{\lstinline!Y!} in the
  tail of \([O, \text{object}]\) or \([Y]\)? No.~We extract
  \passthrough{\lstinline!Y!}:
  \[L(A) = [A, X, Y] + \text{merge}([O, \text{object}], [O, \text{object}])\]
\item
  Evaluate head: \passthrough{\lstinline!O!}. Is
  \passthrough{\lstinline!O!} in the tail of any remaining list?
  No.~Extract \passthrough{\lstinline!O!}:
  \[L(A) = [A, X, Y, O] + \text{merge}([\text{object}], [\text{object}])\]
\item
  Extract \passthrough{\lstinline!object!}:
  \[L(A) = [A, X, Y, O, \text{object}]\]
\end{enumerate}

\hypertarget{the-descriptor-protocol-attribute-lookup-chain}{%
\subsection{4.2 The Descriptor Protocol \& Attribute Lookup
Chain}\label{the-descriptor-protocol-attribute-lookup-chain}}

CPython routes all attribute lookups
(\passthrough{\lstinline!obj.name!}) through the default C function
\passthrough{\lstinline!object\_\_\_getattribute\_\_()!} (which
corresponds to \passthrough{\lstinline!object.\_\_getattribute\_\_!}).

\hypertarget{data-vs.-non-data-descriptors}{%
\subsubsection{Data vs.~Non-Data
Descriptors:}\label{data-vs.-non-data-descriptors}}

\begin{itemize}
\tightlist
\item
  \textbf{Data Descriptor}: Implements both
  \passthrough{\lstinline!\_\_get\_\_!} AND
  \passthrough{\lstinline!\_\_set\_\_!} (or
  \passthrough{\lstinline!\_\_delete\_\_!}).
\item
  \textbf{Non-Data Descriptor}: Implements only
  \passthrough{\lstinline!\_\_get\_\_!} (e.g., standard methods and
  functions).
\end{itemize}

\hypertarget{the-complete-attribute-lookup-resolution-chain}{%
\subsubsection{The Complete Attribute Lookup Resolution
Chain:}\label{the-complete-attribute-lookup-resolution-chain}}

When evaluating \passthrough{\lstinline!obj.name!}, CPython follows this
strict precedence order: 1. \textbf{Data Descriptor}: Checks the class
MRO for a matching data descriptor. If found, returns
\passthrough{\lstinline!Descriptor.\_\_get\_\_(descriptor, obj, Class)!}.
2. \textbf{Instance Dict}: Checks the instance's dictionary
\passthrough{\lstinline!obj.\_\_dict\_\_!} directly. If the key exists,
returns it. 3. \textbf{Non-Data Descriptor / Class Variable}: Checks the
class MRO for a matching non-data descriptor or standard class variable.
If a non-data descriptor is found, returns
\passthrough{\lstinline!Descriptor.\_\_get\_\_(descriptor, obj, Class)!}.
If a class variable is found, returns its raw value. 4. \textbf{Fallback
(Class Dict)}: If still unresolved, checks the class namespace
dictionary for a plain value. 5. \textbf{Fallback Method
(\passthrough{\lstinline!\_\_getattr\_\_!})}: If no attribute is found,
throws \passthrough{\lstinline!AttributeError!} unless the class
implements \passthrough{\lstinline!\_\_getattr\_\_!}, which is then
invoked as a fallback.

\begin{lstlisting}
Attribute Lookup Flowchart:
[Lookup obj.name]
       |
       v
[Class MRO has Data Descriptor?] --Yes--> [Call Descriptor.__get__]
       | No
       v
[Instance dict obj.__dict__ has key?] --Yes--> [Return value from dict]
       | No
       v
[Class MRO has Non-Data Descriptor?] --Yes--> [Call Descriptor.__get__]
       | No
       v
[Class MRO has plain class variable?] --Yes--> [Return class variable]
       | No
       v
[Class defines __getattr__?] --Yes--> [Call __getattr__]
       | No
       +-----------------------------> [Raise AttributeError]
\end{lstlisting}

\hypertarget{method-binding-mechanics}{%
\subsubsection{Method Binding
Mechanics:}\label{method-binding-mechanics}}

Python methods are plain functions attached to a class dictionary.
Because functions implement only \passthrough{\lstinline!\_\_get\_\_!},
they are \textbf{non-data descriptors}. * \textbf{Bound Methods}: When
we call \passthrough{\lstinline!obj.method()!},
\passthrough{\lstinline!method!} is loaded via
\passthrough{\lstinline!\_\_get\_\_!}. The function's
\passthrough{\lstinline!\_\_get\_\_!} method wraps the function and the
instance together, returning a temporary
\passthrough{\lstinline!PyMethod!} object where
\passthrough{\lstinline!self!} is bound to the first argument:
\passthrough{\lstinline!obj.method!} is equivalent to
\passthrough{\lstinline!Class.method.\_\_get\_\_(obj, Class)!}. *
\textbf{Unbound Functions}: If called directly on the class
(\passthrough{\lstinline!Class.method!}),
\passthrough{\lstinline!\_\_get\_\_!} is invoked with
\passthrough{\lstinline!None!} as the instance, returning the raw
function object itself.

\begin{lstlisting}[language=Python]
class Demo:
    def method(self):
        pass

d = Demo()
# Bound method check
print("Bound Method:", d.method)
print("Is bound method object?", type(d.method).__name__)

# Verification of descriptors conversion
bound_manually = Demo.method.__get__(d, Demo)
print("Bound manually matches?", bound_manually == d.method) # True
\end{lstlisting}

\hypertarget{instance-lifecycle-allocation-vs.-initialization}{%
\subsection{4.3 Instance Lifecycle: Allocation
vs.~Initialization}\label{instance-lifecycle-allocation-vs.-initialization}}

CPython splits object creation into two phases, controlled by different
slots: 1. \textbf{Allocation (\passthrough{\lstinline!\_\_new\_\_!})}: *
Maps to the \passthrough{\lstinline!tp\_new!} C slot. * This is a static
method responsible for allocating memory on the heap (using
\passthrough{\lstinline!PyType\_GenericNew()!}) and initializing the
object's \passthrough{\lstinline!PyObject!} headers
(\passthrough{\lstinline!ob\_refcnt!} and
\passthrough{\lstinline!ob\_type!}). * Must return a new instance of the
class (or subclass). 2. \textbf{Initialization
(\passthrough{\lstinline!\_\_init\_\_!})}: * Maps to the
\passthrough{\lstinline!tp\_init!} C slot. * This is an instance method
called immediately after \passthrough{\lstinline!\_\_new\_\_!} returns a
valid instance. It populates the instance dictionary
\passthrough{\lstinline!\_\_dict\_\_!} with variables.

\begin{center}\rule{0.5\linewidth}{0.5pt}\end{center}
